%% =====================================================================
%%  CHAPTER: INTRODUCTION AND BACKGROUND            (budget 3 pages)
%%  chapters/ch1-intro.tex
%% =====================================================================
\pcChapter{intro}{3}{מבוא ורקע}
\begin{pcguide}[מבנה משפך: מהרחב אל הצר]
המבוא פועל כמשפך (איור~\ref{fig:funnel}): מתחילים ברמת המאקרו, יורדים להקשר היישומי, מציגים את המוטיבציה
ומסיימים במבנה המשך הדוח. מסלול הכתיבה: סעיף~\ref{sec:macro} (כ-\latin{1} עמ'), סעיף~\ref{sec:applied} (כ-\latin{0.75}),
סעיף~\ref{sec:motiv} (כ-\latin{0.75}), סעיף~\ref{sec:structure} (כ-\latin{0.5}).
\end{pcguide}

\section{מגמות מאקרו בתחום}\label{sec:macro}
מעבר עובדים מנוסים מתפקידם עלול לקטוע את הקשר בין מידע מתועד לבין שיקול הדעת שנרכש תוך ביצוע העבודה.
סקירה שיטתית על חילופי דורות וניהול ידע מזהה את חשיבותם של תהליכי העברה שיטתיים, תיעוד, למידה חברתית וחניכה
\cite{igoa2024transfer}, ומחקרים עדכניים בתעשייה מצביעים על כך שמודלי שפה יכולים לסייע באיתור ידע סמוי ובשימורו
בתנאי שתהליך הקליטה מובנה ומבוקר \cite{finkel2026unlocking,zuin2025tacit}. בתעשייה עתירת ידע הבעיה אינה רק איתור
קובץ מתאים, אלא גם הבנת ההקשר שבו מידע מסוים תקף, גרסאותיו, החריגים וההשלכות של החלטה שגויה.

התפתחות מערכות \latin{Retrieval-Augmented Generation (RAG)} מאפשרת להפיק תשובות הנשענות על קטעי מקור שאוחזרו ממאגר,
במקום להסתמך רק על הידע הכללי של מודל שפה \cite{lewis2020retrieval}. מסגרות הערכה כגון \latin{RAGAs}, \latin{ARES}
ו-\latin{BERGEN} מספקות כלים לבדיקת רלוונטיות המקורות ונאמנות התשובה \cite{es2024ragas,saad2024ares,rau2024bergen},
וספרות מאוחרת יותר מזהה קשיים מיוחדים במסמכים ארוכים, בטבלאות, בשינויי גרסאות ובשאילתות של משתמשים
\cite{tao2025treerag,strich2026t2,zheng2025knowshift}. במקביל, תופעת ה“הזיה” של מודלי שפה
\cite{ji2023survey} ונקודות הכשל המוכרות בהנדסת מערכות \latin{RAG} \cite{barnett2024seven} מחייבות להפריד בין היכולת לחולל
טקסט לבין נכונותו המקצועית: בסביבה הנדסית פלט שאין לו עיגון ברור במקור, או שאינו תואם לגרסת הנוהל המותרת, לא יועבר
אוטומטית לעובד. לכן הפתרון המוצע יכלול שרשרת עקיבות, ביקורת אנושית ומנגנון הימנעות מתשובה בהיעדר מקור.

\section{ההקשר היישומי}\label{sec:applied}
התעשייה האווירית לישראל \latin{(IAI)} משמשת ארגון הייחוס לפיילוט המוצע. צורכי הארגון שהוצגו לקבוצה הם צמצום אובדן
ידע כאשר עובדים פורשים או עוזבים, קיצור תהליך החפיפה של עובדים חדשים, ושימור מקרים חריגים שנלמדו מניסיון אישי.
לא הועברו לצוות נתוני כוח אדם, מדדי חפיפה, מסמכים פנימיים או הרשאות גישה, ולכן אין בדוח טענה לגבי ממדי הבעיה בפועל
בארגון.

במימוש ארגוני עתידי יוזנו למאגר הידע מסמכים הנדסיים ונהלים שאושרו לעיבוד, ולאחר מכן ייקלטו תיאורי מקרה מעובדים
ותיקים באמצעות ראיונות מובנים. כדי לבצע ניסוי אקדמי ללא חשיפה למידע מסווג, ייבנה תחילה קורפוס חלופי מחומרים
פומביים, ובראשם מדריך ההכשרה הכללי לטכנאי תחזוקת תעופה של רשות התעופה הפדרלית האמריקאית \latin{(FAA)} \cite{faa2023amt}.

\section{מוטיבציה הנדסית וניהולית}\label{sec:motiv}
הערך הניהולי המבוקש הוא מעבר ממסלול חפיפה התלוי בזמינות מומחה לתהליך שניתן לחזור עליו, למדוד אותו ולשפרו.
חלק מן הניסיון מתבטא בזיהוי חזותי של מצב, בבחירת סדר הפעולות ובהתמודדות עם מידע חלקי; לשם כך ייבחן גם תרגול
\latin{3D} ולא רק קריאת מסמך. הערך ההנדסי הוא היכולת לשמר קשר בין שאלת עובד, מקור מדויק, גרסה מאושרת, החלטה
מקצועית, יחידת הדרכה ומבחן כשירות. הפיילוט נועד לבדוק הבטחה זו במדדים מדידים, ולא להניח שמהירות ההפקה מעידה
כשלעצמה על שיפור איכות.

הפרויקט אינו מתכנן הפעלה אוטומטית של מערכות ייצור, הקצאת סמכויות מקצועיות או הפצת מידע מסווג מחוץ לארגון.
התרגול ה\threeDadj{} ייבנה כמודול פיילוט מוגדר: סצנת עבודה מדומה אחת, זיהוי רכיבים, הסתעפויות החלטה ורישום פעולות.
תוקף התוכן יישען על \latin{RAG} ועל אישור מומחה, ותרומת ההמחשה תיבדק בנפרד מול תרגיל \mbox{דו-ממדי}. מחזור ההדרכה המיועד:
עובד ותיק יתאר מקרה מול סצנה כללית, מנהל ידע יאמת את הלקח מול מסמך תקף, מנוע \latin{RAG} יפיק תרחיש ושאלות, מומחה
יאשר, ועובד חדש יתקדם מלמידה ומבחן לתרגול אינטראקטיבי. איור~\ref{fig:funnel} מסכם את המעבר מן המגמות הרחבות אל
הפרויקט עצמו.

\begin{figure}[!htbp]
\centering
\pcFig{funnel}{%
\begin{pctikz}[x=1cm,y=1cm,pcfit=0.86]
  \foreach \i/\w/\t/\c in {0/14/מגמות מאקרו: חילופי דורות ומערכות \latin{RAG}/pcDark, 1/11/הקשר יישומי: חפיפה בתעשייה האווירית/pcPrimary, 2/8/מוטיבציה: חפיפה מדידה ומעוגנת/pcGood, 3/5/הפרויקט: \latin{RAG} + \latin{3D}/pcWarn}{
    \pgfmathsetmacro{\wb}{\w-3}
    \fill[\c!85] ({-\w/2},{-\i*1.05}) -- ({\w/2},{-\i*1.05}) -- ({\wb/2},{-\i*1.05-0.95}) -- ({-\wb/2},{-\i*1.05-0.95}) -- cycle;
    \node[text=white,font=\bfseries\sffamily] at (0,{-\i*1.05-0.48}) {\tn{\t}};
  }
\end{pctikz}}
\caption{משפך המבוא: ממגמות המאקרו אל פרויקט \latin{OrbiTech} (תרשים מבנה; ללא נתונים)}\label{fig:funnel}
\end{figure}

\section{מבנה הדוח}\label{sec:structure}
הדוח מסודר כך שהקורא הבכיר יוכל למצוא בו את מבוקשו בתוך דקות:
\begin{itemize}
  \item פרק~\ref{ch:problem}: הבעיה, שאלת המחקר, ההשערות ומפרט הדרישות.
  \ifdefstring{\SolutionPlacement}{afterproblem}{\item פרק~\ref{ch:solution}: \textbf{הפתרון המוצע והנתונים} --- מה אנו מציעים לבנות ועם איזה מידע.}{}
  \item פרק~\ref{ch:lit}: סקירת ספרות ומאמר הדגל.
  \item פרק~\ref{ch:market}: חקר שוק, מוצרים קיימים ובנצ'מרק.
  \item פרק~\ref{ch:alt}: חלופות תכן, בחירת הפתרון והחדשנות.
  \ifdefstring{\SolutionPlacement}{afteralternatives}{\item פרק~\ref{ch:solution}: \textbf{הפתרון המוצע והנתונים} --- מה אנו מציעים לבנות ועם איזה מידע.}{}
  \item פרק~\ref{ch:plan}: תוכנית עבודה, גאנט ונתיב קריטי.
  \item פרק~\ref{ch:risk}: ניתוח סיכונים ותוכנית הפחתה.
  \item פרק~\ref{ch:summary}: סיכום היתכנות ופעולות מיידיות; לאחריו רשימת המקורות.
\end{itemize}
\pcChapterEnd{intro}
